\documentclass[11pt]{my_memo}
\usepackage{booktabs,tabularx}
\title{Source and Diffuse-LyC Photon-Energy Sampling: Code Comparison}
\author{MoCHII source-code review}\date{2026-07-29}
\begin{document}\maketitle
\section{Scope}
This report is based on direct source inspection of MOCASSIN 2.02.73.2, \texttt{Wood\_Codes/ionize}, CMacIonize-master, \texttt{TORUS/torus/src}, and MoCHII.  For equal-energy packets the source photon CDF is normally proportional to $L_\nu/(h\nu)$.  A free-bound LyC CDF also follows the Milne relation and includes the inverse photoionization cross section.
\section{Comparison}
\begin{tabularx}{\textwidth}{p{22mm}XXX}\toprule
Code&Source energy&Diffuse LyC energy&Result\\\midrule
MoCHII&Continuous input-spectrum inverse CDF; sampled energy retained in continuous mode.&Temperature-indexed Milne CDF per ground continuum, inverted once after interpolating in $\log T$; separate He I cascade branches.&Exact transport energy; exact ground-continuum shape.  He I cascade still approximate.\\
MOCASSIN&Inverse CDF on global frequency grid.&Cell-emissivity CDF on the same grid.&Flexible tabulated emissivity, but grid-valued packets.\\
Wood ionize&Log-interpolated Planck/atmosphere CDF.&Temperature-indexed, cross-section-weighted Milne CDF.&Explicit H/He free-bound continuum shape.\\
CMacIonize&Numerical spectrum-specific CDF.&Dedicated H and He LyC spectra selected by reemission branches.&Modern event-by-event physical reemission.\\
TORUS&Tabulated wavelength CDF, linear in selected interval.&No active HII diffuse-LyC reemission path found.&Generic continuum sampler only.\\\bottomrule
\end{tabularx}
\section{MoCHII}
\texttt{src/energy\_\allowbreak sampler\_\allowbreak mod.f90} constructs the source CDF and \texttt{src/main.f90} launches packets.  In \texttt{ion\_\allowbreak energy\_\allowbreak mode} $=$ \texttt{'continuous'}, \texttt{assign\_\allowbreak ion\_\allowbreak packet\_\allowbreak energy} preserves the sampled energy; \texttt{ion\_\allowbreak packet\_\allowbreak mod.f90} and \texttt{ion\_\allowbreak score\_\allowbreak mod.f90} evaluate H, He, and metal physics at that energy.  The ionizing-band index is diagnostic only.
\texttt{src/diffuse\_mod.f90} builds a leaf--channel CDF each iteration.  Its H II$\to$H I(1s), He II$\to$He I(1$^1$S), and He III$\to$He II(1s) channels use $\alpha_1=\alpha_A-\alpha_B$ and, since 2026-07-29, sample the Milne photon spectrum $p(E)\propto E^2\sigma_{\rm pi}(E)e^{-(E-E_{th})/kT_e}$ through \texttt{src/milne\_\allowbreak recomb\_\allowbreak spectrum\_\allowbreak mod.f90}, selected by \texttt{par\%diffuse\_\allowbreak energy\_\allowbreak model} $=$ \texttt{'milne'} (the default).  The tables are built once in $x=(E-E_{th})/kT$, interpolated in $\log T$, and inverted once; both pseudorandom and Sobol launches call the same routine, and the Sobol dimension assignment is unchanged.  Because packets are equal-energy and their count comes from the channel's energy luminosity, they are drawn from $E\,n(E)$ rather than $n(E)$, which makes the represented photon rate exactly $n_en_{\rm ion}\alpha_1V$; the luminosity still uses the number-weighted mean $\langle E\rangle_n$.  Start-up reports $\alpha_1$ recovered from the same spectral integral against $\alpha_A-\alpha_B$ without rescaling either.  \texttt{exponential} ($E=E_{th}+kT_e[-\ln u]$) is retained for reproducing earlier runs and \texttt{threshold} as a sensitivity bracket; neither is a physical replacement.  The optional He I cascade still selects 19.82-eV and 21.22-eV branches by energy luminosity and samples its ionizing two-photon branch uniformly from 13.598--20.62 eV, and is now the crudest remaining piece.
\section{MOCASSIN}
In \texttt{source/photon\_mod.f90}, \texttt{newPhotonPacket} calls \texttt{getNu2(inSpectrumProbDen,...)} for stellar photons.  \texttt{getNu} and \texttt{getNu2} invert cumulative arrays and return an integer frequency pointer.  \texttt{continuum\_mod.f90} builds the continuum representation and \texttt{setLdiffuse} accumulates diffuse luminosity.  Diffuse continuum packets use the same pointer/CDF mechanism.  Thus arbitrary free-bound emissivities can be tabulated, but packet frequency is a global-grid value rather than a continuous in-bin variate.
\section{Wood ionize}
\texttt{planckset.f} prepares the source CDF; \texttt{nuset.f} locates a uniform deviate in \texttt{bbcdf} and logarithmically interpolates \texttt{lognu}.  \texttt{probinit.f} explicitly constructs photon CDFs, $\int j_\nu/\nu\,d\nu$.  At every \texttt{Tarray} point it calls \texttt{phfit2} and uses $j_{HI}\propto\nu^3\sigma_{HI}(\nu)\exp[-157919.667(\nu-1)/T]$, with an analogous He I expression above 1.81 Ryd.  It integrates $j_\nu/\nu$ to \texttt{pH} and \texttt{pHe}.  \texttt{nuemit.f} inverts the temperature-indexed H/He CDFs and the He two-photon CDF; \texttt{mcionize.f} selects the reemission branch using \texttt{probset.f}.  This is a Milne-grid sampler, though its output is a grid frequency.
\section{CMacIonize}
\texttt{PlanckPhotonSourceSpectrum.cpp} tabulates $\nu^2/(e^{h\nu/kT}-1)$ over 1--4 Ryd, integrates a CDF, and interpolates in logarithmic CDF--frequency space in \texttt{get\_random\_frequency}.  \texttt{PhysicalDiffuseReemissionHandler.cpp} chooses absorption with $n_i\sigma_i$ weights and applies temperature-dependent reemission probabilities.  It invokes \texttt{HydrogenLymanContinuumSpectrum} and \texttt{HeliumLymanContinuumSpectrum}, and separately treats the 19.8-eV and 56\% ionizing He two-photon branches.  These dedicated spectra receive cross-section objects, implementing the same physical design as Wood ionize.
\section{TORUS}
\texttt{spectrum\_mod.f90:getWavelength} draws $r$, locates it in \texttt{spectrum\%prob}, and linearly interpolates adjacent tabulated wavelengths; \texttt{photon\_mod.f90:initPhoton} uses it.  Searches of active \texttt{torus/src} found no Wood/CMacIonize-like H--He reemission handler, temperature-indexed free-bound LyC CDF, or HII recombination-ionizing-packet emission call.  This is an absence of exposed HII diffuse-LyC physics, not evidence for a threshold approximation.
\section{Implication}
MoCHII already evaluates transport cross sections at each sampled source energy, unlike frequency-pointer approaches.  Its ground-state diffuse continuum, formerly the distinct limitation, now follows the Milne relation with the same threshold-aware cross section transport uses, interpolating the CDF in temperature before a single inversion --- the prescription \S\ref{sec:accurate} argues for, and one step beyond both Wood (nearest temperature row, grid frequency) and CMacIonize (interpolation between frequencies drawn from two temperature rows).  The measured outcome confirms the prediction made here: the converged HII40 calculation moves the He I front by only $+0.6\times10^{-3}$ pc, closing $0.3\%$ of the Cloudy residual.  The exact sampler was physically required and it did not remove the residual.  Two leads remain: the He I excited cascade, still using fixed AGN3 branching and a flat two-photon surrogate, which alone moves that front by $0.10$ pc; and a $2$--$3.4\%$ internal disagreement between $\alpha_1$ from the He I Milne integral and the fitted $\alpha_A-\alpha_B$, in exactly the continuum whose front is discrepant.
\section{Which method is physically most accurate?}\label{sec:accurate}
The most accurate practical Monte-Carlo prescription is \emph{not} a fixed energy, an exponential tail, or merely a very fine arbitrary energy grid.  For each recombining ion, electron temperature, and (where required) density, construct the photon-number PDF from the level-resolved free--bound emissivity obtained by detailed balance,
\[p_i(\nu\mid T_e)\propto {j^{\rm fb}_{\nu,i}(T_e)\over h\nu},\qquad \nu\ge\nu_{0,i},\]
using the same resolved photoionization cross section in the Milne relation that is used to absorb the packet.  The emissivity should include the relevant free--bound Gaunt factor and should be integrated with accurate quadrature; its CDF and its mean photon energy must be derived from that same PDF.  Branching to recombination levels and subsequent He I/He II cascades must be selected from level-resolved effective coefficients, with the full two-photon spectral shape rather than a flat surrogate.  This prescription is exact subject to the adopted atomic data and the assumption of a Maxwellian electron distribution.
Wood ionize and CMacIonize are closest among the reviewed codes: both use $\nu^3\sigma(\nu)\exp[-h(\nu-\nu_0)/kT]$ and temperature tables.  CMacIonize improves the source CDF interpolation, but its LyC implementation interpolates frequencies selected from two neighbouring temperature CDFs rather than directly in a temperature-interpolated CDF; Wood uses the nearest temperature row.  Therefore neither is the final reference implementation.  MoCHII now does interpolate the \emph{CDF} in temperature, invert once, sample a continuous energy, and carry that same energy into the existing exact transport cross sections.  The remaining approximation in its diffuse field is no longer the ground-continuum shape but the He I excited cascade, where a level-resolved calculation and a physical two-photon spectrum are still owed.
\section{Why can Wood ionize be fast?}
The accuracy of the diffuse-LyC \emph{spectral shape} is not the dominant cost.  Wood ionize precomputes its small $T$--$\nu$ Milne CDF arrays in \texttt{probinit.f}; during transport, \texttt{nuemit.f} only draws a uniform deviate and performs a table search.  It does not numerically integrate the Milne relation or evaluate a new continuum spectrum for every emitted packet.

More importantly, Wood uses event-by-event reemission.  Following absorption, \texttt{mcionize.f} selects absorption species and reemission channel; a diffuse packet is created only when that event branches into an ionizing reemission.  MoCHII instead rebuilds an explicit volume-emissivity CDF every iteration and transports a separate diffuse population whose size is $L_{\rm diffuse}/L_{\rm packet}$.  In HII40 this is commonly $10^7$--$2\times10^7$ diffuse packets per iteration, in addition to direct packets.

Wood also confines this benchmark physics largely to H and He and samples a modest frequency grid.  MoCHII additionally pays for AMR leaf traversal, continuous-energy H/He/metal path scoring, metal ionization cascades and cooling, thermal equilibrium iterations, and HDF5 diagnostics.  These choices can be scientifically necessary, but they dominate runtime relative to one inverse-CDF operation.

Consequently, adding a precomputed Milne CDF to MoCHII should have negligible direct sampling cost.  It should not be expected by itself to make MoCHII as fast as Wood: the major cost is explicit diffuse-packet transport and the broader MoCHII physics.  The correct comparison is therefore limited: Wood is closest for the H/He ground-continuum spectral PDF, but its nearest-temperature, grid-frequency sampling is still less rigorous than a MoCHII CDF/PDF-temperature interpolation followed by continuous-energy inversion.
\end{document}
